\subsection{Annuities}


\content{
        Formulas so far: \\
        \fbox{\begin{minipage}{\textwidth}
        For these formulas, we have the following: 
        \begin{itemize}
        \item $P_0$ is the pricipal. 
        \item $r$ is the interest rate. 
        \item $t$ is the number of time intervals (specified by problem) which
        have passed. 
        \item $A$ is the amount paid back.
        \item $I$ is the interest.
        \end{itemize}
        \begin{center}
        \begin{tabular}{|c|c|}
        \hline
        Simple Interest & \begin{minipage}[t!]{1.5in} \begin{eqnarray*} I=P_0r \\
                A = P_0(1+r)\end{eqnarray*}\end{minipage} \\
        \hline
        Simple Interest over Time & \begin{minipage}{1.5in} \begin{eqnarray*}
        I=P_0rt \\ A = P_0(1+rt) \end{eqnarray*}\end{minipage}\\
        \hline
        \end{tabular}
        \end{center}
        \end{minipage}}
        \fbox{\begin{minipage}{\textwidth}
        For these formulas, we have the following: 
        \begin{itemize}
        \item $P_0$ is the pricipal. 
        \item $r$ is the annual interest rate.
        \item $t$ is the number of years which have passed. 
        \item $A$ is the total amount.
        \item $k$ is the number of times interest is compounded per year. 
        \end{itemize}
        \begin{center}
        \begin{tabular}{|c|c|}
        \hline
        Compound Interest & \begin{minipage}{1.5in}$$ A =
        P_0\left(1+\frac{r}{k}\right)^{kt} $$\end{minipage}\\
        \hline
        \end{tabular}
        \end{center}
        \end{minipage}}
}{% presentation
}{% nopresentation
}{% comments
}

\content{
        \begin{example} You will need \$40,000 for your childs education in 18
        years. You open a new savings account with a 4\% APR compounded
        quarterly. How much will you need to invest into the account now to have
        enough money when your child is 18?
        \end{example}
}{% presentation
\vspace{2in}
}{% nopresentation
\vspace{2in}
}{% comments
}

\content{
        \newpage
        \begin{definition} (Savings Annuity) Placing a small amount in an
        account each month to slowly build up savings instead of placing a large
        sum to begin with.
        \end{definition}
}{% presentation
\vspace{4in}
}{% nopresentation
\vspace{4in}
}{% comments
Go over derivation: 
\begin{enumerate}
\item DEFINE $r=APR$, $k$ - number of compounds per year. $d$ the deposit
amount.
\item Write recursive formula for $P_m$, where $m$ is the number of compounds
that have occured. 
\item Expand, and write two of these
\item multiply second one by $(1+r/k)$ then subtract to simplify. 
\end{enumerate}
Final formula 
$P_m = d \frac{\left(\left(1+\frac{r}{k}\right)^m-1\right)}{\frac{r}{k}} $
one last step, $m=tk$, where $t$ is the number of years that have passed. 

\textbf{Note:} in this formula we assume that we make a deposit once every time
the interest is compounded, so if there is a problem where the number of
compounds per year is not stated, just use the number of deposits per year for
$k$. 
}

\content{
        \begin{example} Suppose you create an individual
        retirement account (IRA) with a 6\% APR which is compounded monthly.
        Each month you deposit \$100 into the account. How much will be in the
        account in 20 years? 
        \end{example}
}{% presentation
\vspace{3in}
}{% nopresentation
\vspace{3in}
}{% comments
}

\content{
        \begin{example} You want to have \$200,000 in your account when you 
        retire in 30 years. Your retirement account earns 8\% interest. 
        How much do you need to deposit each month to meet your
        retirement goal?
        \end{example}
}{% presentation
\vspace{3in}
}{% nopresentation
\vspace{3in}
}{% comments
}

\content{
        \begin{example} You open a new bank account with a 3\% APR, compounded
        daily, and each day you deposit \$10 into it. How much will you have in
        the account in 10 years? 
        \end{example}
}{% presentation
\vspace{3in}
}{% nopresentation
\vspace{3in}
}{% comments
}

\content{
        \begin{example} You want to have \$10,000 dollars saved in a special
        savings account for a trip in 3 years. The account has a 2.5\% APR
        compounded twice a month. If you plan on depositing money into the
        account twice a month, how much will you need to deposit to reach your
        goal? 
        \end{example}
}{% presentation
\vspace{3in}
}{% nopresentation
\vspace{3in}
}{% comments
}


